% Provenance (never printed): see sources/outline_sources.md (Year 2: eccentricity, obliquity and
% precession; insolation computed from spherical astronomy; the Quaternary cycles; carbon dioxide
% and feedbacks). Orbital elements and tabulated insolation of Berger and Loutre (1991) via NOAA
% NCEI (data/derived/bachelor-2-orbit91, fetch_data Book 4 block); periods after Laskar et al.
% (2004) and the NASA Milankovitch page; history after the NASA Earth Observatory feature (2000);
% terminations after Barker and Knorr (2021); the mid-Pleistocene transition after Clark et al.
% (2006) and Koehler and van de Wal (2020). Figures from figdata/bachelor-2/orbital.py; values of
% the exercises and problem computed in the Book 4 scratch calculator calc/ch25.py.
\chapter{Forçage orbital et âges glaciaires}\label{ch:b2:orbital-forcing}

Au cours du dernier million d'années, la forme de l'orbite terrestre a modifié
la lumière reçue du Soleil par toute la planète en une année de moins de deux pour mille~:
entre \num{340.00} et \qty{340.54}{W.m^{-2}} au sommet de l'atmosphère. Pendant le même
million d'années, l'insolation journalière moyenne de mi-juillet à \ang{65}~N a varié entre
environ 388 et \qty{500}{W.m^{-2}}, soit plus d'un quart. Les âges glaciaires ont suivi le second
nombre, et non le premier. Ce chapitre calcule les deux à partir des éléments de l'orbite
terrestre, et demande comment le rayonnement solaire, à une latitude et pendant une saison donnée, peut rythmer la croissance et la fonte des calottes glaciaires qui recouvrent les continents.
% ledger: data:b4-orbital, berger91:readme

\begin{recall}
\textbf{Des volumes de physique}~: l'ellipse, son demi-grand axe $a$ et son excentricité $e$, les lois
de Kepler. \textbf{Du volume de première année}~: l'écliptique et l'obliquité $\varepsilon$, la
déclinaison $\delta$ du Soleil, avec $\sin\delta = \sin\varepsilon\sin\lambda$ ($\lambda$ la longitude vraie du
Soleil, comptée le long de l'écliptique à partir de l'équinoxe de printemps), l'année tropique, l'insolation
journalière moyenne
\[
  \bar Q = \frac{S_0}{\pi}\left(\frac{a}{r}\right)^{2}\bigl(h_0\sin\varphi\sin\delta +
  \cos\varphi\cos\delta\sin h_0\bigr),\qquad \cos h_0 = -\tan\varphi\tan\delta ,
\]
et sa moyenne globale et annuelle, $S_0/(4\sqrt{1 - e^2})$. \textbf{Des volumes scolaires}~: le
cycle glaciaire et le dernier maximum glaciaire~; le forçage radiatif. \textbf{Des chapitres précédents}~: la
\omterm{def:b2:precambrian-earth:ice-albedo}{rétroaction glace--albédo} (\cref{def:b2:precambrian-earth:ice-albedo})~; l'empilement benthique, l'enregistrement EPICA
de Dôme C et les \omterm{def:b2:climate-proxies:age-model}{modèles d'âge} (\cref{ch:b2:climate-proxies}).
\end{recall}
% ledger: earth:TSI, earth:obliquity, earth:ecc

\section{L'orbite et ses trois éléments}\label{sec:b2:orbital-forcing:elements}

L'orbite terrestre est une ellipse dont le Soleil occupe l'un des foyers. Les planètes, surtout Jupiter et Saturne,
la déforment lentement et la font tourner dans son plan ; l'attraction du Soleil et de la Lune sur le
renflement équatorial fait osciller l'axe de rotation comme une toupie. Trois grandeurs décrivent ce qui
compte pour le rayonnement solaire : l'allongement de l'ellipse, l'inclinaison de l'axe, et l'endroit de
l'orbite où tombent les saisons.
% ledger: nasa:milank-causes

\begin{definition}[Longitude du périhélie~; précession climatique]\label{def:b2:orbital-forcing:elements}
La \emph{longitude du périhélie}\index{longitude du périhélie} $\varpi$ est l'angle, mesuré dans le plan de l'orbite,
entre la direction de l'équinoxe de printemps et la direction du périhélie vue depuis le Soleil. Elle est comptée
à partir de l'équinoxe mobile~: elle varie donc avec le mouvement de l'axe et avec la rotation de
l'orbite. La \emph{précession
climatique}\index{précession climatique} est la quantité $e\sin\varpi$, appelée aussi indice de
précession.
\end{definition}
% ledger: berger91:readme

\begin{omfigure}
\centering
\begin{tikzpicture}[scale=0.82, font=\scriptsize, >=Stealth]
\draw[thick, rotate=101.37] (-0.72,0) ellipse (2.4 and 2.2895);
\draw[densely dashed, black!60] (-0.331,1.647) -- (0.615,-3.059);
\draw[->, black!70] (0,0) -- (3.9,0);
\node[anchor=south, align=center] at (3.75,0.06) {$\gamma$, vers l'équinoxe\\de printemps};
\draw[->, omDef] (0.75,0) arc[start angle=0, end angle=101.37, radius=0.75];
\node[omDef] at (0.78,0.62) {$\varpi$};
\fill[ocRed] (0,0) circle (0.13);
\node[anchor=north west] at (0.05,-0.08) {Soleil};
\foreach \x/\y in {-2.062/0, 0/-3.094, 2.321/0, 0/1.688} {\fill[ocBlue] (\x,\y) circle (0.08);}
\node[anchor=east, align=right] at (-2.15,0) {équinoxe\\de mars};
\node[anchor=north west, align=left] at (2.38,-0.08) {équinoxe\\de septembre};
\node[anchor=west] at (0.1,1.92) {solstice de décembre};
\node[anchor=east] at (-0.12,-3.25) {solstice de juin};
\fill[black] (-0.331,1.647) circle (0.06);
\node[anchor=east] at (-0.42,1.72) {périhélie};
\fill[black] (0.615,-3.059) circle (0.06);
\node[anchor=west] at (0.7,-3.15) {aphélie};
\draw[->, thick, black!60] (-1.9,-1.75) arc[start angle=220, end angle=245, radius=2.6];
\end{tikzpicture}
\par\smallskip{\small Orbite terrestre vue du côté nord, avec une excentricité exagérée à 0.3 et la
longitude actuelle du périhélie, $\varpi = \ang{101.4}$. La Terre parcourt son orbite dans le sens inverse des aiguilles d'une montre (flèche)~; les points marquent sa position aux
équinoxes et aux solstices. Le périhélie tombe quelques jours après le solstice de décembre,
l'aphélie quelques jours après le solstice de juin. Schéma.}
\end{omfigure}
% ledger: berger91:rows, nasa:milank-detail

Vue depuis la Terre, le Soleil se déplace le long de l'écliptique, et sa longitude vraie $\lambda$ est l'angle
compté à partir de l'équinoxe de printemps~: $\lambda = \ang{90}$ au solstice de juin, \ang{270} au solstice de décembre.
La Terre est au périhélie lorsque le Soleil se trouve dans la direction opposée à celle du périhélie,
soit $\lambda = \varpi + \ang{180}$~; aujourd'hui $\varpi = \ang{101.37}$ et le périhélie tombe à
$\lambda = \ang{281.37}$, environ onze jours après le solstice de décembre, vers le 3 janvier.
% ledger: berger91:rows, nasa:milank-detail

\begin{definition}[Paramètres de Milankovitch]\label{def:b2:orbital-forcing:milankovitch}
Les \emph{paramètres de Milankovitch}\index{paramètres de Milankovitch} sont les trois grandeurs orbitales
qui fixent la répartition de l'ensoleillement sur la Terre selon les saisons et les latitudes~: l'excentricité $e$,
l'obliquité $\varepsilon$ et la \omterm{def:b2:orbital-forcing:elements}{précession climatique} $e\sin\varpi$. La \emph{théorie de
Milankovitch}\index{théorie de Milankovitch} affirme que leurs variations lentes rythment la croissance
et la fonte des calottes glaciaires du Quaternaire, par l'intermédiaire de l'insolation estivale aux
hautes latitudes nord.
\end{definition}
% ledger: nasa:milank-detail, eo00:milankovitch

Chacun de ces paramètres varie selon son propre rythme. L'excentricité a des périodes principales voisines de \num{95000},
\num{99000}, \num{124000} et \num{131000} ans, qui se combinent en un cycle d'environ
\num{100000} ans, et une période de \num{405000} ans, la plus stable de toutes~; sur le dernier
million d'années, elle a varié de \num{0.0046} à \num{0.0565}, et aujourd'hui, à \num{0.0172}, elle diminue lentement.
L'obliquité varie avec une période principale de \num{40978} ans, entre \ang{22.1} et
\ang{24.5} sur le dernier million d'années~; à \ang{23.44}, elle diminue, après son dernier maximum
vers \num{10000} ans avant aujourd'hui. L'axe oscille en environ \num{25770} ans, tandis que
l'orbite tourne en environ \num{112000} ans~; ensemble, ils font passer le périhélie par les saisons, avec des
périodes principales voisines de \num{19000}, \num{22400} et \num{23700} ans.
% ledger: la04:eccentricity, la04:obliquity, la04:precession, nasa:milank-detail, data:b4-orbital

\begin{omfigure}
\centering
% figdata/bachelor-2/orbital.py part elements
\begin{tikzpicture}
\begin{groupplot}[group style={group size=1 by 3, vertical sep=0.3cm, xlabels at=edge bottom,
  xticklabels at=edge bottom}, width=0.96\linewidth, height=3.1cm, xmin=0, xmax=1000,
  xtick={0,100,...,1000}, tick label style={font=\scriptsize}, label style={font=\small},
  scaled y ticks=false, yticklabel style={/pgf/number format/fixed,
  /pgf/number format/precision=2}]
\nextgroupplot[ylabel={$e$}, ymin=0, ymax=0.085, ytick={0,0.02,0.04,0.06}]
\addplot[omDef, thin] table[x=age_ka, y=ecc] {figdata/out/bachelor-2-orbital-elements.dat};
\node[font=\scriptsize, anchor=north west] at (rel axis cs:0.01,0.97) {a\enspace excentricité};
\nextgroupplot[ylabel={$\varepsilon$ (\unit{\degree})}, ymin=21.9, ymax=25.4, ytick={22,23,24}]
\addplot[ocBlue, thin] table[x=age_ka, y=obl] {figdata/out/bachelor-2-orbital-elements.dat};
\node[font=\scriptsize, anchor=north west] at (rel axis cs:0.01,0.97) {b\enspace obliquité};
\nextgroupplot[ylabel={$e\sin\varpi$}, ymin=-0.07, ymax=0.12, ytick={-0.05,0,0.05},
  xlabel={âge (milliers d'années avant 1950)}]
\addplot[ocRed, thin] table[x=age_ka, y=prec] {figdata/out/bachelor-2-orbital-elements.dat};
\node[font=\scriptsize, anchor=north west] at (rel axis cs:0.01,0.97) {c\enspace précession
climatique};
\end{groupplot}
\end{tikzpicture}
\par\smallskip{\small Les trois \omterm{def:b2:orbital-forcing:milankovitch}{paramètres de Milankovitch} sur le dernier million d'années~: (a) l'excentricité~;
(b) l'obliquité~; (c) la \omterm{def:b2:orbital-forcing:elements}{précession climatique}, dont l'enveloppe est l'excentricité.
Données~: Berger et Loutre (1991), NOAA NCEI Paleoclimatology, consultées le 23 octobre 2026.}
\end{omfigure}
% ledger: data:b4-orbit91

La \omterm{def:b2:orbital-forcing:elements}{précession climatique} oscille vite, mais seulement dans la mesure où l'excentricité le permet~: lorsque l'orbite est
presque circulaire, comme il y a \num{400000} ans, peu importe où tombe le périhélie. La figure le montre
clairement~: l'enveloppe du panneau (c) est le panneau (a).
% ledger: data:b4-orbit91

\section{Rayonnement solaire issu de l'orbite}\label{sec:b2:orbital-forcing:insolation}

L'insolation journalière rappelée plus haut demande un ingrédient de plus~: la distance du Soleil
à un jour donné.

\begin{proposition}[Distance et longitude du Soleil]\label{prop:b2:orbital-forcing:distance}
Lorsque la longitude vraie du Soleil vaut $\lambda$, la distance Terre--Soleil vaut
\[
  \frac{r}{a} = \frac{1 - e^2}{1 - e\cos(\lambda - \varpi)} .
\]
Au solstice de juin ($\lambda = \ang{90}$), au premier ordre en $e$,
$(a/r)^2 \approx 1 - 2e\sin\varpi$~; au solstice de décembre, $(a/r)^2 \approx 1 + 2e\sin\varpi$.
\end{proposition}

\begin{proof}
Sur une ellipse dont le Soleil occupe un foyer, $r/a = (1 - e^2)/(1 + e\cos v)$, où l'anomalie vraie $v$
est l'angle, vu du Soleil, entre le périhélie et la Terre. La Terre est vue du Soleil dans la direction
opposée à celle sous laquelle le Soleil est vu depuis la Terre, si bien que sa longitude vue du Soleil vaut
$\lambda + \ang{180}$, et $v = \lambda + \ang{180} - \varpi$, $\cos v = -\cos(\lambda - \varpi)$.
Pour $\lambda = \ang{90}$, $\cos(\lambda - \varpi) = \sin\varpi$ et $(a/r)^2 = (1 -
e\sin\varpi)^2/(1 - e^2)^2 = 1 - 2e\sin\varpi + O(e^2)$~; pour $\lambda = \ang{270}$, le signe change.
\end{proof}

La \omterm{def:b2:orbital-forcing:elements}{précession climatique} est donc le bon indice : elle mesure à quel point la Terre est plus proche
du Soleil, ou plus éloignée, au solstice de juin. Lorsque $e\sin\varpi$ est grand et positif, comme aujourd'hui et
bien davantage il y a \num{116000} ans, l'été boréal tombe près de l'aphélie et il est faible~; lorsqu'il est grand et négatif,
comme il y a environ \num{10000} ans, l'été boréal tombe près du périhélie.
% ledger: berger91:rows

\begin{example}[Le Soleil de juin à \ang{65}~N, aujourd'hui]\label{ex:b2:orbital-forcing:today}
Aujourd'hui $e = 0.017236$, $\varpi = \ang{101.37}$, $\varepsilon = \ang{23.446}$, et les tables utilisent
$S_0 = \qty{1360}{W.m^{-2}}$. Au solstice de juin $\delta = \varepsilon$~; à $\varphi =
\ang{65}$, $\cos h_0 = -\tan\ang{65}\tan\ang{23.446} = -0.9301$, donc $h_0 = \ang{158.44} =
\qty{2.7654}{rad}$~: le Soleil est au-dessus de l'horizon pendant $2 \times 158.44/360 \times 24 = 21.1$~heures. Le facteur entre parenthèses de la formule d'insolation vaut
$2.7654 \times 0.9063 \times 0.3979 + 0.4226 \times 0.9174 \times 0.3674 =
1.1397$, et à la distance moyenne $\bar Q = 432.90 \times 1.1397 = \qty{493.4}{W.m^{-2}}$. Le
facteur de distance vaut $(a/r)^2 = 0.9671$, donc $\bar Q = \qty{477.1}{W.m^{-2}}$.
\end{example}
% ledger: berger91:rows, berger91:readme, data:b4-orbital

La seconde loi de Kepler ajoute une nuance. Près de l'aphélie, la Terre se déplace plus lentement : la saison qui tombe là
est plus faible, mais dure plus longtemps. Aujourd'hui, les étés boréaux durent environ 4.5 jours de plus que les hivers boréaux,
et les printemps environ trois jours de plus que les automnes.
% ledger: nasa:milank-detail

\begin{proposition}[Les moyennes annuelles ne dépendent pas de la précession]\label{prop:b2:orbital-forcing:annual}
L'énergie reçue par unité de surface au sommet de l'atmosphère entre deux longitudes données du
Soleil ne dépend pas de $\varpi$, et l'insolation annuelle moyenne à toute latitude est celle d'une orbite circulaire de rayon $a$
multipliée par $1/\sqrt{1 - e^2}$. Au pôle, la moyenne annuelle vaut
\[
  \bar Q_{\text{pôle}} = \frac{S_0\sin\varepsilon}{\pi\sqrt{1 - e^2}} ,
\]
qui croît avec l'obliquité.
\end{proposition}

\begin{proof}
Par la seconde loi de Kepler, le rayon balaie une aire à la vitesse constante $\pi ab/T$, et l'aire balayée
pendant que la Terre tourne de $\mathrm d\lambda$ vaut $\frac12 r^2\,\mathrm d\lambda$, donc $\mathrm dt =
T r^2\,\mathrm d\lambda/(2\pi ab)$, avec $b = a\sqrt{1 - e^2}$. L'insolation à une latitude donnée et pour une longitude donnée du
Soleil vaut $(a/r)^2$ fois sa valeur $Q_a(\lambda)$ à la distance $a$, si bien que
l'énergie reçue est $\int Q_a(\lambda)(a/r)^2\,\mathrm dt = \frac{T}{2\pi\sqrt{1 - e^2}}\int
Q_a(\lambda)\,\mathrm d\lambda$~: la distance s'est éliminée, et avec elle $\varpi$. Sur toute l'année,
$\lambda$ parcourt \ang{360} et la moyenne vaut $\frac{1}{2\pi\sqrt{1-e^2}}\int_0^{2\pi}
Q_a\,\mathrm d\lambda$. Au pôle Nord $Q_a = S_0\sin\delta = S_0\sin\varepsilon\sin\lambda$
pour $0 < \lambda < \pi$ et 0 sinon~; l'intégrale vaut $2S_0\sin\varepsilon$.
\end{proof}
% ledger: earth:TSI

La précession déplace le rayonnement solaire d'une saison à l'autre, sans jamais le faire entrer ni sortir d'une année~; l'excentricité
modifie le total annuel d'un facteur $1/\sqrt{1 - e^2}$, de \qty{0.16}{\percent} au plus sur le dernier million d'années~;
l'obliquité déplace le rayonnement entre les latitudes, en élevant la moyenne annuelle aux pôles et en l'abaissant
légèrement à l'équateur. Seuls le régime saisonnier et la répartition
latitudinale changent beaucoup, et c'est ce que montre la figure ci-dessous.
% ledger: data:b4-orbital, nasa:milank-detail

\begin{omfigure}
\centering
% figdata/bachelor-2/orbital.py part lat
\begin{tikzpicture}
\begin{groupplot}[group style={group size=2 by 1, horizontal sep=1.6cm},
  tick label style={font=\scriptsize}, label style={font=\small}, height=5.4cm,
  width=0.48\linewidth, xmin=-90, xmax=90, xtick={-90,-60,...,90}, xlabel={latitude
  (\unit{\degree})}, ylabel={$\bar Q$ (\unit{W.m^{-2}})}]
\nextgroupplot[ymin=0, ymax=650, title={a\enspace solstice de juin}, title style={font=\scriptsize},
  legend style={font=\scriptsize, at={(0.03,0.97)}, anchor=north west, draw=none},
  legend cell align=left]
\addplot[black, thick] table[x=lat, y=jun_0] {figdata/out/bachelor-2-orbital-lat.dat};
\addlegendentry{aujourd'hui}
\addplot[ocRed, thick, densely dashed] table[x=lat, y=jun_10] {figdata/out/bachelor-2-orbital-lat.dat};
\addlegendentry{il y a \num{10000} ans}
\nextgroupplot[ymin=100, ymax=500, title={b\enspace moyenne annuelle}, title style={font=\scriptsize},
  legend style={font=\scriptsize, at={(0.5,0.03)}, anchor=south, draw=none},
  legend cell align=left]
\addplot[ocBlue, thick] table[x=lat, y=ann_lo] {figdata/out/bachelor-2-orbital-lat.dat};
\addlegendentry{$\varepsilon = \ang{22.1}$}
\addplot[ocRed, thick, densely dashed] table[x=lat, y=ann_hi] {figdata/out/bachelor-2-orbital-lat.dat};
\addlegendentry{$\varepsilon = \ang{24.5}$}
\end{groupplot}
\end{tikzpicture}
\par\smallskip{\small (a) Insolation journalière moyenne au solstice de juin, aujourd'hui et il y a \num{10000} ans, lorsque
l'été boréal tombait près du périhélie et que l'obliquité était plus forte~: environ
\qty{50}{W.m^{-2}} de plus dans les hautes latitudes nord. (b) Insolation annuelle moyenne pour les deux obliquités
extrêmes du dernier million d'années, avec l'excentricité actuelle~: une inclinaison plus grande envoie plus de
rayonnement vers les pôles et un peu moins vers les tropiques. Calculé à partir des éléments orbitaux
avec $S_0 = \qty{1360}{W.m^{-2}}$.}
\end{omfigure}
% ledger: data:b4-orbital, data:b4-orbit91

\section{Calcul de la courbe d'insolation}\label{sec:b2:orbital-forcing:curve}

La courbe d'insolation construite par Milankovitch est aujourd'hui un calcul de routine. Les éléments orbitaux
du passé se calculent en intégrant les mouvements des planètes~; tabulés tous les mille ans,
ils servent de données d'entrée. Pour chaque âge, la méthode ci-dessous donne l'insolation à toute latitude et pour toute saison.

\begin{method}[Insolation du passé]\label{met:b2:orbital-forcing:curve}
\begin{enumerate}
  \item Lire $e$, $\varpi$ et $\varepsilon$ pour l'âge voulu.
  \item Choisir la saison par la longitude vraie du Soleil $\lambda$~: \ang{90} pour le solstice de juin~;
        le mi-juillet des tables correspond à $\lambda = \ang{120}$.
  \item Calculer $\delta$ à partir de $\sin\delta = \sin\varepsilon\sin\lambda$, puis l'angle horaire du
        coucher du Soleil $h_0$ à la latitude $\varphi$ (jour ou nuit polaire si $|\tan\varphi\tan\delta| > 1$).
  \item Calculer le facteur de distance $(a/r)^2$ avec la \cref{prop:b2:orbital-forcing:distance}.
  \item Multiplier~: $\bar Q = (S_0/\pi)(a/r)^2(h_0\sin\varphi\sin\delta +
        \cos\varphi\cos\delta\sin h_0)$.
\end{enumerate}
\end{method}
% ledger: berger91:readme, data:b4-orbital

Appliquée aux 1001 âges du dernier million d'années, la méthode reproduit l'insolation de mi-juillet tabulée
à \ang{65}~N à \qty{0.012}{W.m^{-2}} près, soit l'arrondi des tables, et reproduit exactement les valeurs de juin
tabulées pour aujourd'hui. Les tables placent l'origine des temps en 1950.
% ledger: data:b4-orbital, berger91:readme

\begin{omfigure}
\centering
% figdata/bachelor-2/orbital.py part jul65
\begin{tikzpicture}
\begin{groupplot}[group style={group size=1 by 2, vertical sep=0.3cm, xlabels at=edge bottom,
  xticklabels at=edge bottom}, width=0.96\linewidth, xmin=0, xmax=400, xtick={0,50,...,400},
  tick label style={font=\scriptsize}, label style={font=\small}]
\nextgroupplot[height=5.0cm, ymin=370, ymax=590, ylabel={$\bar Q$ à \ang{65}~N
  (\unit{W.m^{-2}})}, legend style={font=\scriptsize, at={(0.5,1.02)}, anchor=south, draw=none,
  legend columns=3, /tikz/every even column/.append style={column sep=6pt}},
  legend cell align=left]
\addplot[black!45, line width=2.2pt] table[x=age_ka, y=tab] {figdata/out/bachelor-2-orbital-jul65.dat};
\addlegendentry{mi-juillet, valeurs tabulées}
\addplot[ocRed, thin] table[x=age_ka, y=calc] {figdata/out/bachelor-2-orbital-jul65.dat};
\addlegendentry{mi-juillet, calculé}
\addplot[ocBlue, thin, densely dashed] table[x=age_ka, y=jun65] {figdata/out/bachelor-2-orbital-jul65.dat};
\addlegendentry{solstice de juin, calculé}
\nextgroupplot[height=2.8cm, ymin=-0.02, ymax=0.02, ytick={-0.01,0,0.01},
  scaled y ticks=false, yticklabel style={/pgf/number format/fixed}, ylabel={résidu},
  xlabel={âge (milliers d'années avant 1950)}]
\addplot[only marks, mark=*, mark size=0.35pt, black] table[x=age_ka, y=resid]
  {figdata/out/bachelor-2-orbital-jul65.dat};
\end{groupplot}
\end{tikzpicture}
\par\smallskip{\small En haut : insolation journalière moyenne à \ang{65}~N sur les \num{400000} dernières années,
de mi-juillet telle que tabulée par Berger et Loutre (gris, épais) et telle que calculée par la méthode ci-dessus (rouge),
et au solstice de juin (bleu, pointillé). En bas : valeurs calculées moins valeurs tabulées de mi-juillet, en
\unit{W.m^{-2}}.}
\end{omfigure}
% ledger: data:b4-orbit91, data:b4-orbital

Deux traits ressortent. La courbe est dominée par la \omterm{def:b2:orbital-forcing:elements}{précession climatique}, avec des pics espacés d'environ
\num{20000} ans, dont la hauteur croît et décroît avec l'excentricité~; l'obliquité ajoute une modulation plus lente.
Et les variations sont grandes~: sur les \num{800000} dernières années, la valeur de mi-juillet
a varié de 388 à \qty{497}{W.m^{-2}}.
% ledger: data:b4-orbital

\section{Théorie de Milankovitch et cycles glaciaires}\label{sec:b2:orbital-forcing:theory}

\begin{history}[Ensoleillement à soixante-cinq degrés nord]
Milutin Milankovitch, né en 1879, a bâti une théorie mathématique du climat à partir des variations saisonnières et latitudinales de l'ensoleillement. En combinant les trois variations orbitales,
il calcula l'insolation des différentes latitudes et les températures de surface qui en découlent pour les \num{600000} années précédant 1800.
Sur la suggestion du climatologue Köppen, il retint comme grandeur décisive l'insolation estivale à \ang{65}~N~: les grandes calottes glaciaires s'étaient développées
près de cette latitude, et des étés frais fondent moins la neige de l'hiver, si bien que la neige s'accumule
d'une année sur l'autre. Il tenait l'obliquité pour le plus important des trois paramètres. Pendant une cinquantaine
d'années, la théorie fut largement ignorée~; Milankovitch mourut en 1958. En 1976, Hays, Imbrie et Shackleton trouvèrent
dans les sédiments des grands fonds,
pour les derniers \num{450000} ans, des pics spectraux de la variabilité climatique
vers \num{23000}, \num{42000} et \num{100000} ans, les périodes de la précession, de l'obliquité
et de l'excentricité, et conclurent que les variations de l'orbite terrestre
sont la cause fondamentale des âges glaciaires du Quaternaire.
\end{history}
% ledger: eo00:milankovitch, nasa:milank-detail, hays76:periods

La figure ci-dessous reprend leur test sur les enregistrements du \cref{ch:b2:climate-proxies}. Elle compare les spectres d'amplitude, c'est-à-dire la force de chaque période,
de l'insolation de juillet à \ang{65}~N et de l'empilement benthique,
dans deux fenêtres de \num{700000} ans.

\begin{omfigure}
\centering
% figdata/bachelor-2/orbital.py part spectrum
\begin{tikzpicture}
\begin{axis}[width=0.9\linewidth, height=5.6cm, xmin=10, xmax=150, ymin=0, ymax=1.08,
  xlabel={période (milliers d'années)}, ylabel={amplitude relative},
  xtick={19,23,41,100,150}, tick label style={font=\scriptsize}, label style={font=\small},
  legend style={font=\scriptsize, at={(0.5,1.02)}, anchor=south, draw=none, legend columns=3,
  /tikz/every even column/.append style={column sep=6pt}}, legend cell align=left]
\addplot[black!50, thick] table[x=period, y=amp_ins] {figdata/out/bachelor-2-orbital-spectrum.dat};
\addlegendentry{insolation à \ang{65}~N, 0--\qty{700}{ka}}
\addplot[ocRed, thick] table[x=period, y=amp_young] {figdata/out/bachelor-2-orbital-spectrum.dat};
\addlegendentry{LR04, 0--\qty{700}{ka}}
\addplot[ocBlue, thick, densely dashed] table[x=period, y=amp_old] {figdata/out/bachelor-2-orbital-spectrum.dat};
\addlegendentry{LR04, \numrange{1250}{1950}~ka}
\end{axis}
\end{tikzpicture}
\par\smallskip{\small Spectres d'amplitude, chacun normalisé à son propre maximum, de l'insolation de mi-juillet à \ang{65}~N sur les
derniers \num{700000} ans et de l'empilement benthique LR04 sur la même fenêtre et sur \num{1250000} à \num{1950000} ans avant aujourd'hui
(tendance retirée, fenêtre atténuée). L'insolation a sa puissance aux périodes de la précession, voisines de 19 à 24 mille ans,
et à la période de l'obliquité~; le volume de glace récent culmine vers
\num{100000} ans, là où l'insolation n'a presque rien~; l'ancien volume de
glace culmine à \num{41000} ans.}
\end{omfigure}
% ledger: data:b4-orbital, data:lr04, hays76:periods

\subsection*{Glaciaires, interglaciaires et terminaisons}

\begin{definition}[Période glaciaire~; interglaciaire~; terminaison glaciaire]\label{def:b2:orbital-forcing:glacials}
Dans l'enregistrement du Quaternaire, une \emph{période glaciaire}\index{période glaciaire} est un intervalle où de grandes calottes glaciaires recouvrent
les continents de l'hémisphère Nord, ce que traduisent un $\delta^{18}$O benthique élevé et un niveau marin bas~;
un \emph{interglaciaire}\index{interglaciaire} est un intervalle où la glace continentale de l'hémisphère Nord est
faible hors du Groenland, comme aujourd'hui. Une \emph{terminaison
glaciaire}\index{terminaison glaciaire} est l'une des grandes transitions qui mettent fin à un cycle glaciaire,
du plein glaciaire aux conditions interglaciaires, beaucoup plus vite que la croissance qui l'a précédée.
\end{definition}
% ledger: kohler20:mpt, barker21:terminations, ipcc:ref-periods

Sur le dernier demi-million d'années, la période dominante du climat est voisine de \num{100000} ans,
bien que le forçage à cette période, qui provient de l'excentricité seule, soit négligeable
à côté de celui de l'obliquité et de la précession. Les cycles sont aussi dissymétriques : le $\delta^{18}$O benthique
enregistre une accumulation de glace par paliers pendant environ \num{90000} ans, puis une décroissance en environ
\num{10000} ans, en dents de scie dont les fronts descendants sont les terminaisons. Le dernier \omterm{def:b2:orbital-forcing:glacials}{interglaciaire}, d'environ
\num{129000} à \num{116000} ans, s'acheva lorsque l'insolation estivale à \ang{65}~N tomba à l'une de ses valeurs
les plus basses des derniers \num{800000} ans : la valeur de juin atteignit son minimum,
\qty{438.6}{W.m^{-2}}, il y a \num{116000} ans, et la valeur de mi-juillet son second minimum,
\qty{392.2}{W.m^{-2}}, il y a \num{114000} ans. Le dernier maximum glaciaire survint
environ \num{95000} ans plus tard.
% ledger: barker21:terminations, ipcc:ref-periods, data:b4-orbital, data:b4-orbital-extremes

Avec quelle rapidité une calotte glaciaire peut-elle suivre son forçage~? Le modèle le plus simple traite le volume de glace $V$ comme
tendant vers le volume $V_{\mathrm{eq}}$ que soutiendrait l'insolation actuelle, avec un
temps de réponse $\tau$~:
\[
  \frac{\mathrm dV}{\mathrm dt} = \frac{V_{\mathrm{eq}}(t) - V}{\tau} .
\]

\begin{proposition}[Une réponse retardée]\label{prop:b2:orbital-forcing:lag}
Pour un forçage $V_{\mathrm{eq}} = A\sin\omega t$, la solution du modèle à retard tend vers
\[
  V = \frac{A}{\sqrt{1 + \omega^2\tau^2}}\,\sin\bigl(\omega(t - t_{\ell})\bigr),\qquad
  t_{\ell} = \frac{\arctan(\omega\tau)}{\omega} :
\]
la glace suit presque entièrement un forçage lent et seulement faiblement un forçage rapide, toujours en retard. La réponse
à une somme de sinusoïdes est la somme des réponses~: le modèle ne crée aucune
période qui manque à son forçage.
\end{proposition}

\begin{proof}
Essayons $V = B\sin\omega t + C\cos\omega t$~: en substituant, $\omega B\tau = -C$ et $-\omega C\tau = A -
B$, donc $B = A/(1 + \omega^2\tau^2)$, $C = -\omega\tau A/(1 + \omega^2\tau^2)$, une amplitude
$\sqrt{B^2 + C^2} = A/\sqrt{1 + \omega^2\tau^2}$ et un retard de phase $\arctan(\omega\tau)$. Toute autre
solution diffère de celle-ci d'un multiple de $\mathrm e^{-t/\tau}$, qui s'éteint. L'équation
est linéaire, donc les solutions s'additionnent.
\end{proof}

\begin{omfigure}
\centering
% figdata/bachelor-2/orbital.py part lag
\begin{tikzpicture}
\begin{axis}[width=0.96\linewidth, height=5.4cm, xmin=0, xmax=400, ymin=-3.4, ymax=4.2,
  xtick={0,50,...,400}, xlabel={âge (milliers d'années avant 1950)},
  ylabel={normalisé}, tick label style={font=\scriptsize}, label style={font=\small},
  legend style={font=\scriptsize, at={(0.5,1.02)}, anchor=south, draw=none, legend columns=3,
  /tikz/every even column/.append style={column sep=6pt}}, legend cell align=left]
\addplot[black!35, thin] table[x=age_ka, y=forcing] {figdata/out/bachelor-2-orbital-lag.dat};
\addlegendentry{moins l'insolation à \ang{65}~N}
\addplot[ocRed, thick] table[x=age_ka, y=model] {figdata/out/bachelor-2-orbital-lag.dat};
\addlegendentry{modèle à retard, $\tau = \qty{10}{ka}$}
\addplot[black, thick] table[x=age_ka, y=lr04] {figdata/out/bachelor-2-orbital-lag.dat};
\addlegendentry{LR04 (plus de glace)}
\end{axis}
\end{tikzpicture}
\par\smallskip{\small Le modèle à retard, piloté par l'insolation de mi-juillet à \ang{65}~N (moins d'ensoleillement, plus de glace),
comparé à l'empilement benthique, les trois séries étant normalisées (moyenne 0, écart type 1)~; des valeurs positives signifient plus de glace.
Le modèle suit les oscillations de précession de l'empilement mais n'a aucun de ses grands cycles en dents de scie.
Le temps de réponse de \qty{10}{ka} est un choix du modèle.}
\end{omfigure}
% ledger: data:b4-orbital, data:lr04

Sur les derniers \num{800000} ans, le modèle à retard est corrélé à l'empilement avec $r = 0.42$, et l'insolation seule
avec seulement $0.23$~: le retard aide, mais une calotte glaciaire linéaire ne peut pas produire des cycles de
\num{100000} ans à partir d'un forçage qui n'en a aucun. Les terminaisons doivent faire intervenir un comportement non linéaire qui amplifie
une faible poussée de l'insolation. Les candidats sont des rétroactions qui agissent sur l'ensemble du climat~:
la \omterm{def:b2:precambrian-earth:ice-albedo}{rétroaction glace--albédo} de la \cref{def:b2:precambrian-earth:ice-albedo}, puisqu'une planète couverte de neige et de glace
réfléchit davantage le rayonnement solaire~; et le dioxyde de carbone, qui, dans les carottes de glace antarctiques,
augmente à l'échelle de millénaires chaque fois que la circulation de retournement de l'Atlantique s'affaiblit,
comme c'est le cas lorsque l'eau de fonte des calottes glaciaires se déverse dans l'Atlantique Nord. Au-delà d'un seuil
de dioxyde de carbone, les conditions \omterm{def:b2:orbital-forcing:glacials}{interglaciaires} peuvent se stabiliser.
% ledger: data:b4-orbital, barker21:terminations, eo00:milankovitch

\begin{definition}[Transition du Pléistocène moyen]\label{def:b2:orbital-forcing:mpt}
La \emph{transition du Pléistocène moyen}\index{transition du Pléistocène moyen} est le passage, d'environ \num{1250000} à \num{700000} ans,
de cycles glaciaires d'environ \num{41000} ans, de faible amplitude, à des cycles quasi périodiques d'environ \num{100000} ans,
de grande amplitude, avec des calottes glaciaires plus épaisses, sans aucun changement correspondant
du forçage orbital.
\end{definition}
% ledger: clark06:mpt, nasa:milank-detail, kohler20:mpt

Avant la transition, les calottes glaciaires suivaient l'obliquité, comme le montre le spectre en pointillé de la figure
ci-dessus~; un décompte trouve une \omterm{def:b2:orbital-forcing:glacials}{terminaison glaciaire} dans 67, 88 et \qty{52}{\percent} des cycles d'obliquité
du Quaternaire ancien, moyen et récent. La théorie orbitale seule ne peut pas expliquer le changement~; les candidats sont
la dynamique des calottes glaciaires elles-mêmes, peut-être après l'enlèvement du régolithe meuble sur lequel elles glissaient,
et une diminution progressive du dioxyde de carbone pendant les minima glaciaires. Aucune explication n'est encore
généralement admise.
% ledger: kohler20:mpt, nasa:milank-detail, clark06:mpt

\section{Le dioxyde de carbone dans les cycles glaciaires}\label{sec:b2:orbital-forcing:co2}

Les carottes de glace antarctiques montrent que le dioxyde de carbone suit les cycles glaciaires : entre 188 et
\qty{194}{ppm} au dernier maximum glaciaire, entre 266 et \qty{282}{ppm} au dernier \omterm{def:b2:orbital-forcing:glacials}{interglaciaire}~;
la série composite du \cref{ch:b2:climate-proxies} donne \qty{189.7}{ppm} pour 19 à
\qty{23}{ka}, et \qty{279.4}{ppm} sur les deux mille ans précédant l'industrialisation. Son forçage radiatif
croît à peu près comme le logarithme de la concentration, \qty{3.93}{W.m^{-2}} par doublement.
% ledger: ipcc:fig2-34, data:b4-proxies-co2, ipcc:log-co2, ipcc:erf-2xco2

\begin{example}[L'orbite face au dioxyde de carbone]\label{ex:b2:orbital-forcing:co2}
De \qty{189.7}{ppm} à \qty{279.4}{ppm}, le forçage vaut $3.93\log_2(279.4/189.7) = 3.93 \times
0.559 = \qty{2.2}{W.m^{-2}}$, sur tout le globe et pendant toute l'année. L'orbite, sur le dernier million d'années, a modifié
la moyenne globale et annuelle de l'insolation d'au plus $340.54 - 340.00 =
\qty{0.54}{W.m^{-2}}$, mais l'insolation de mi-juillet à \ang{65}~N d'environ \qty{109}{W.m^{-2}}.
\end{example}
% ledger: data:b4-proxies-co2, ipcc:erf-2xco2, ipcc:log-co2, data:b4-orbital

L'orbite est un métronome, non un moteur. Elle redistribue le rayonnement solaire entre les saisons et les latitudes
et fixe le calendrier~; les calottes glaciaires, par leur albédo et leur lenteur, et le dioxyde de carbone, par son forçage
global, transforment un signal régional en un changement global du climat.
% ledger: hays76:periods, barker21:terminations, eo00:milankovitch

\section*{Exercices}\label{sec:b2:orbital-forcing:exercises}

Données des exercices, sauf indication contraire~: $S_0 = \qty{1360}{W.m^{-2}}$ (la valeur des
tables)~; une année tropique de \num{365.2422} jours~; les éléments orbitaux des tables, dont ceux d'aujourd'hui
sont $e = 0.017236$, $\varpi = \ang{101.37}$, $\varepsilon = \ang{23.446}$.
% ledger: berger91:readme, berger91:rows

\begin{exercise}[$\star$]\label{exo:b2:orbital-forcing:1}
Prendre $e = 0.0167$. (a) Calculer le rapport de la distance à l'aphélie $a(1 + e)$ à la distance au périhélie
$a(1 - e)$. (b) Exprimer leur différence en pourcentage de $a$. (c) De quel facteur le rayonnement solaire au périhélie
est-il supérieur à celui de l'aphélie~?
\end{exercise}
% ledger: earth:ecc, nasa:milank-detail

\begin{exercise}[$\star$]\label{exo:b2:orbital-forcing:2}
Les tables donnent $e = 0.017236$, $\varpi = \ang{101.37}$ aujourd'hui~; $e = 0.019971$, $\varpi =
\ang{293.86}$ il y a \num{10000} ans~; $e = 0.019249$, $\varpi = \ang{359.99}$ il y a \num{6000} ans.
(a) Calculer la \omterm{def:b2:orbital-forcing:elements}{précession climatique} pour chaque âge. (b) Pour chacun, donner la longitude du Soleil
au périhélie. (c) À quel âge le périhélie était-il le plus proche du solstice de juin, et à quel âge
était-il le plus proche de l'équinoxe de septembre~?
\end{exercise}
% ledger: berger91:rows

\begin{exercise}[$\star$]\label{exo:b2:orbital-forcing:3}
L'axe de rotation précesse avec une période d'environ \num{25770} ans par rapport aux étoiles, tandis que
l'orbite tourne avec une période d'environ \num{112000} ans. Supposer que les deux mouvements déplacent
le périhélie par rapport à l'équinoxe dans le même sens, de sorte que leurs fréquences s'ajoutent. (a) Combiner
les deux fréquences pour trouver la période avec laquelle le périhélie parcourt les saisons. (b) Comparer cette période aux
périodes principales de la \omterm{def:b2:orbital-forcing:elements}{précession climatique},
\num{19.0}, \num{22.4} et \num{23.7} milliers d'années.
\end{exercise}
% ledger: nasa:milank-detail, la04:precession

\begin{exercise}[$\star$]\label{exo:b2:orbital-forcing:4}
Avec la \cref{prop:b2:orbital-forcing:annual} et $e = 0.0167$, calculer l'insolation annuelle moyenne au
pôle Nord pour des obliquités de \ang{22.1} et \ang{24.5}. De quel pourcentage croît-elle entre
les deux~?
\end{exercise}
% ledger: nasa:milank-detail, earth:ecc

\begin{exercise}[$\star\star$]\label{exo:b2:orbital-forcing:5}
(a) Calculer exactement le facteur de distance $(a/r)^2$ au solstice de juin, aujourd'hui et il y a \num{116000}
ans ($e = 0.043988$, $\varpi = \ang{92.71}$), et comparer avec l'approximation $1 -
2e\sin\varpi$. (b) Calculer les deux au solstice de décembre. (c) Pourquoi l'approximation est-elle moins bonne il y a
\num{116000} ans~?
\end{exercise}
% ledger: berger91:rows

\begin{exercise}[$\star\star$]\label{exo:b2:orbital-forcing:6}
Reprendre pas à pas l'\cref{ex:b2:orbital-forcing:today}~: l'angle horaire du coucher du Soleil à \ang{65}~N
au solstice de juin, la durée du jour, l'insolation à la distance
moyenne, puis l'insolation d'aujourd'hui.
\end{exercise}
% ledger: berger91:rows

\begin{exercise}[$\star\star$]\label{exo:b2:orbital-forcing:7}
(a) En partant de $\mathrm dt = T(r/a)^2\,\mathrm d\lambda/(2\pi\sqrt{1 - e^2})$ et de $(r/a)^2 \approx 1 +
2e\cos(\lambda - \varpi)$, montrer que la demi-année allant de l'équinoxe de mars à celui de septembre dure environ
$T/2 + (2T/\pi)\,e\sin\varpi$. (b) Calculer de combien de jours elle dépasse l'autre demi-année
aujourd'hui. (c) Comparer avec les 4.5 jours dont les étés boréaux dépassent les hivers, et les trois jours dont
les printemps dépassent les automnes.
\end{exercise}
% ledger: nasa:milank-detail, berger91:rows

\begin{exercise}[$\star\star$]\label{exo:b2:orbital-forcing:8}
Au cours du dernier million d'années, l'excentricité a varié de \num{0.0046} à \num{0.0565}. (a) Calculer
l'insolation moyenne globale et annuelle $S_0/(4\sqrt{1 - e^2})$ aux deux extrêmes. (b) De combien, en
\unit{W.m^{-2}} et en pourcentage, l'orbite a-t-elle modifié cette valeur~?
(c) Comparer avec l'intervalle de l'insolation de mi-juillet à \ang{65}~N, de 388 à \qty{497}{W.m^{-2}}.
\end{exercise}
% ledger: data:b4-orbital

\begin{exercise}[$\star\star$]\label{exo:b2:orbital-forcing:9}
Avec le modèle à retard et $\tau = \qty{10}{ka}$, calculer l'amplitude relative et le retard de la masse de glace
pour un forçage sinusoïdal de périodes 23, 41 et \num{100} milliers d'années. Quelle bande la calotte glaciaire
favorise-t-elle, et pourquoi~?
\end{exercise}

\begin{exercise}[$\star\star\star$]\label{exo:b2:orbital-forcing:10}
L'insolation de mi-juillet à \ang{65}~N n'a presque aucune puissance à \num{100000} ans, pourtant le volume de glace
des derniers \num{700000} ans y concentre l'essentiel de sa puissance. (a) Montrer que si le forçage vaut $A_1\sin
\omega_1 t + A_2\sin\omega_2 t$, le modèle à retard ne répond qu'aux fréquences $\omega_1$ et $\omega_2$. (b) Montrer
qu'une réponse contenant le terme $V^2$, par exemple $\mathrm dV/\mathrm dt = (V_{\mathrm{eq}} -
V)/\tau + kV^2$, peut créer les fréquences $\omega_1 - \omega_2$ et $\omega_1 + \omega_2$, en
substituant la solution linéaire dans $V^2$. (c) Deux termes de précession de périodes \num{23.7} et
\num{19.0} milliers d'années battent avec quelle période~?
Que conclure concernant les cycles de \num{100000} ans~?
\end{exercise}
% ledger: la04:precession, barker21:terminations

\begin{exercise}[$\star\star\star$]\label{exo:b2:orbital-forcing:11}
Au pôle Nord, l'énergie reçue par mètre carré pendant le jour polaire vaut $E = \int S_0(a/r)^2
\sin\delta\,\mathrm dt$ sur la demi-année où $\delta > 0$. (a) Montrer, avec la démonstration de
\cref{prop:b2:orbital-forcing:annual}, que $E = S_0T\sin\varepsilon/(\pi\sqrt{1 - e^2})$. (b) Calculer
$E$ en \unit{GJ.m^{-2}} pour $\varepsilon = \ang{22.1}$ et \ang{24.5}, avec $e = 0.0167$ et $T$ =
365.25 jours. (c) La précession n'apparaît pas dans $E$. Pourquoi peut-elle néanmoins modifier la fonte
d'une calotte glaciaire~?
\end{exercise}
% ledger: nasa:milank-detail

\begin{exercise}[$\star\star\star$]\label{exo:b2:orbital-forcing:12}
(a) Calculer le forçage radiatif de la hausse du dioxyde de carbone de \qty{189.7}{ppm} à
\qty{279.4}{ppm}, avec \qty{3.93}{W.m^{-2}} par doublement et une loi logarithmique. (b) Les intervalles
donnés pour le dernier maximum glaciaire et le dernier \omterm{def:b2:orbital-forcing:glacials}{interglaciaire} sont de 188 à \qty{194}{ppm} et de 266 à
\qty{282}{ppm}~: calculer le forçage le plus faible et le plus fort qu'ils permettent. (c) Comparer avec la variation de
l'insolation annuelle globale de l'\cref{exo:b2:orbital-forcing:8}
et avec l'amplitude de l'insolation estivale à \ang{65}~N. Lequel est le déclencheur, et lequel l'amplificateur~?
\end{exercise}
% ledger: data:b4-proxies-co2, ipcc:fig2-34, ipcc:erf-2xco2, ipcc:log-co2, data:b4-orbital

\begin{problem}[{Problème du week-end --- Rayonnement solaire à \ang{65}~N}]\label{pb:b2:orbital-forcing:1}
Le dernier \omterm{def:b2:orbital-forcing:glacials}{interglaciaire} s'est achevé il y a environ \num{116000} ans. Quel ensoleillement l'été à la latitude des grandes calottes glaciaires de l'hémisphère Nord
a-t-il reçu alors, par rapport à aujourd'hui~? Utiliser les données des exercices.
Les tables donnent, il y a \num{116000} ans, $e = 0.0440$, $\varpi = \ang{92.7}$,
$\varepsilon = \ang{22.52}$~; aujourd'hui, arrondies de la même façon, $e = 0.0172$, $\varpi = \ang{101.4}$,
$\varepsilon = \ang{23.45}$.

\medskip\noindent\textbf{Partie I --- L'orbite alors et aujourd'hui.}
\begin{enumerate}
  \item Calculer la \omterm{def:b2:orbital-forcing:elements}{précession climatique} alors et aujourd'hui.
  \item Calculer les distances au périhélie et à l'aphélie, en unités de $a$, alors et aujourd'hui.
  \item À quelle longitude du Soleil la Terre était-elle au périhélie, alors et aujourd'hui~? Près de quel solstice~?
  \item De quel facteur le rayonnement solaire au périhélie était-il supérieur à celui de l'aphélie, alors et aujourd'hui~?
  \item L'obliquité était-elle plus grande ou plus petite qu'aujourd'hui~? Qu'est-ce que cela signifie pour les étés des
        hautes latitudes~?
  \item Calculer la moyenne globale et annuelle de l'insolation alors et aujourd'hui. La différence est-elle significative~?
\end{enumerate}

\medskip\noindent\textbf{Partie II --- La distance au solstice de juin.}
\begin{enumerate}[resume]
  \item Montrer à partir de la \cref{prop:b2:orbital-forcing:distance} qu'au solstice de juin $r/a = (1 -
        e^2)/(1 - e\sin\varpi)$.
  \item Calculer $r/a$ au solstice de juin, alors et aujourd'hui.
  \item Calculer $(a/r)^2$ alors et aujourd'hui.
  \item Comparer avec l'approximation $1 - 2e\sin\varpi$.
  \item De quel pourcentage le rayonnement du solstice était-il plus faible qu'à la distance moyenne, alors et aujourd'hui~?
\end{enumerate}

\medskip\noindent\textbf{Partie III --- Ensoleillement à \ang{65}~N.}
\begin{enumerate}[resume]
  \item Donner la déclinaison du Soleil au solstice de juin, alors et aujourd'hui.
  \item Calculer l'angle horaire du coucher du Soleil à \ang{65}~N ce jour-là, alors et aujourd'hui.
  \item En déduire la durée du jour.
  \item Calculer l'insolation journalière moyenne à la distance moyenne.
  \item Calculer l'insolation journalière moyenne, alors et aujourd'hui.
  \item Les tables donnent, pour le mi-juillet, 397.60 et \qty{426.76}{W.m^{-2}}. Pourquoi les valeurs
        au solstice sont-elles plus élevées~?
  \item Lequel pèse le plus dans la différence entre alors et aujourd'hui, l'obliquité ou la distance~?
\end{enumerate}

\medskip\noindent\textbf{Partie IV --- La réponse de la calotte glaciaire.}
\begin{enumerate}[resume]
  \item Dans le modèle à retard avec $\tau = \qty{10}{ka}$, quelle fraction d'un changement brusque de
        $V_{\mathrm{eq}}$ la calotte glaciaire a-t-elle atteinte après 5, 10 et \qty{20}{ka}~?
  \item L'insolation de juin à \ang{65}~N est passée d'un maximum, survenu il y a environ \num{128000} ans, à
        le minimum survenu il y a \num{116000} ans, soit une demi-période de précession. Pour une sinusoïde de période
        \qty{23}{ka}, quelle fraction du forçage la calotte glaciaire suit-elle, et avec quel retard~?
  \item Pourquoi une calotte glaciaire ne peut-elle pas atteindre, en un cycle de précession, la taille qu'une faible
        insolation estivale soutiendrait~?
  \item La glace a mis environ \num{90000} ans à se former jusqu'au dernier maximum glaciaire, et environ
        \num{10000} ans à fondre. Pourquoi le modèle à retard ne peut-il pas produire cette dissymétrie~?
  \item Le dioxyde de carbone se situait entre 266 et \qty{282}{ppm} au dernier \omterm{def:b2:orbital-forcing:glacials}{interglaciaire}, et entre
        188 et \qty{194}{ppm} au dernier maximum glaciaire. Comment aide-t-il un changement régional de
        l'ensoleillement à devenir un changement global du climat~?
  \item En résumé : quelle était l'insolation journalière moyenne à \ang{65}~N au solstice de juin lorsque le dernier \omterm{def:b2:orbital-forcing:glacials}{interglaciaire}
        s'est achevé, il y a \num{116000} ans~?
\end{enumerate}
\end{problem}
% ledger: berger91:rows, berger91:readme, data:b4-orbital, ipcc:ref-periods, ipcc:fig2-34, barker21:terminations
